\section{谱方法：用光滑性换高精度}
\label{sec:08-Numerical-Analysis/08-Fourier-and-Spectral-Methods/04}

手上是一个周期区间 \([0,2\pi)\) 上的 Poisson 方程 \(-u''=f\)，要求解的精度到 \(10^{-10}\)。取 \(u(x)=e^{\sin x}\) 造一个有精确解的算例，好把误差量出来。

二阶中心差分给出的误差按 \(h^2\) 走。要压到 \(10^{-10}\)，步长得取到 \(10^{-5}\) 量级，网格点数以几十万计，还得解一个几十万阶的三对角系统。换成 Fourier 谱方法：把 \(f\) 做一次 FFT，逐个频率除以 \(k^2\)，再做一次逆变换。取 \(32\) 个模式，误差已经落在 \(10^{-16}\) 附近——不是 \(10^{-10}\)，是舍入水平。三十二个自由度对几十万个。

把右端换一下，让解变成周期三角波那样带一个角点的函数，同一份代码跑出来的数字是另一回事：\(N=32\) 时误差 \(10^{-2}\)，加到 \(N=2048\)，误差 \(1.6\times10^{-4}\)。\(N\) 放大六十四倍，误差降到六十四分之一。这是 \(O(N^{-1})\)，一个彻底的代数阶。

\begin{center}
\begin{tikzpicture}[scale=1.0,font=\small,>=stealth]
  \draw[->,black!55] (-0.15,0) -- (7.6,0) node[right,font=\footnotesize] {\(N\)};
  \draw[->,black!55] (0,0.25) -- (0,-6.3);
  \foreach \yy/\lab in {0/{\(10^{0}\)},-1.75/{\(10^{-5}\)},-3.5/{\(10^{-10}\)},-5.25/{\(10^{-15}\)}}{
    \draw[black!45] (-0.1,\yy) -- (0.1,\yy);
    \node[left,font=\footnotesize,black!70] at (-0.14,\yy) {\lab};
  }
  \foreach \nn in {16,32,48,64}{
    \draw[black!45] ({\nn*0.109375},0.1) -- ({\nn*0.109375},-0.1);
    \node[above,font=\footnotesize,black!70] at ({\nn*0.109375},0.12) {\nn};
  }
  \draw[black!35,dashed] (0,-5.6) -- (7.3,-5.6);
  \node[right,font=\footnotesize,black!55] at (5.0,-5.85) {舍入平台};
  \draw[blue!65!black,thick] (0.219,-0.35) -- (3.28,-5.6);
  \draw[blue!65!black,thick] (3.28,-5.6) -- (7.0,-5.6);
  \node[right,font=\footnotesize,blue!55!black] at (2.15,-3.15) {光滑解：直线};
  \draw[orange!80!black,thick] plot[domain=2:64,samples=60] ({\x*0.109375},{log10(2/\x)*0.35});
  \node[right,font=\footnotesize,orange!70!black] at (4.2,-0.72) {带角点的解：弯曲且趋平};
  \node[font=\footnotesize,black!70] at (3.6,-6.75) {截断误差（纵轴对数刻度）};
\end{tikzpicture}

\emph{同一个求解器、同一个方程，只换了解的光滑性。蓝线每往右走一格就掉一个固定的倍数，直到撞上舍入平台；橙线走到右边界还在同一个数量级里挪动。}
\end{center}

蓝线是直线这件事值得盯一会儿。半对数坐标下的直线意味着 \(\text{误差}\approx C\rho^{N}\)，\(\rho<1\)：每多留一个基函数，误差就乘以一个固定的常数，图中这个常数大约是 \(0.3\)。差分方法给的是 \(h^{p}\)，\(p\) 是格式定死的阶——网格加倍，误差降 \(2^{p}\) 倍，这个倍数永远是 \(2^{p}\)。谱方法给的是复利。橙线则说明，同一个求解器在光滑性不达标时连差分都不如。

\subsection{微分在频率坐标里退化成乘法}

有限差分在每个点上只看邻近几个节点，用一个低次多项式局部拟合再求导。谱方法反过来，把整个未知函数摊到一组全局基上：

\[
u_N(x)=\sum_{k}\hat u_k\phi_k(x),
\]

周期问题取 \(\phi_k(x)=e^{ikx}\)，有限非周期区间取 Chebyshev 多项式。全局基的好处在求导那一步立刻兑现：

\[
u(x)=\sum_k\hat u_ke^{ikx}\ \Longrightarrow\ u'(x)=\sum_k(ik)\hat u_ke^{ikx},\qquad u''(x)=\sum_k(-k^2)\hat u_ke^{ikx}.
\]

微分算子在这组坐标里是对角的。于是 \(-u''=f\) 在每个非零频率上塌缩成一个标量方程 \(k^2\hat u_k=\hat f_k\)，逐模式除一下就完事。整个求解过程是两次 FFT 加一遍除法，成本 \(O(N\log N)\)，这正是\hyperref[sec:08-Numerical-Analysis/08-Fourier-and-Spectral-Methods/03]{``DFT 是变换，FFT 是算法''}那条流水线的直接用法。

\(k=0\) 是个例外，也是第一个容易踩的坑。常数函数落在 \(-u''\) 的零空间里，方程在零频上写成 \(0\cdot\hat u_0=\hat f_0\)：若 \(\hat f_0\ne0\) 则无解，若 \(\hat f_0=0\) 则 \(\hat u_0\) 任意。前者是兼容条件（右端在周期区间上必须均值为零），后者要靠额外指定 \(u\) 的均值来定。代码里忘掉零模的表现是除零或者结果整体漂移一个常数。

方程怎么施加在展开式上，有几种传统做法。Galerkin 要求残差与所有试验基正交，得到的系数方程最干净但非线性项难算；tau 方法保留大部分模态方程，把最高几阶换成边界条件；collocation（也叫 pseudospectral）在一组选定节点上直接令残差为零，做法上是在节点值与系数之间来回变换，工程实现里最常见。三者在光滑问题上收敛速度相当，差别落在边界处理和非线性项的代价上。

\subsection{指数收敛来自系数衰减，与算法是否精巧无关}

谱方法的精度全部由一件事决定：截断掉的尾巴有多小。截断误差就是 \(\sum_{\abs{k}>N/2}\hat u_ke^{ikx}\) 这一截，它的大小直接由系数衰减率给出，算法本身没有任何魔法可加。

\hyperref[sec:08-Numerical-Analysis/08-Fourier-and-Spectral-Methods/01]{``从 Fourier 级数到连续变换''}已经建立了衰减与光滑性的对应：\(f\) 有 \(m\) 阶可积导数则 \(\hat f(k)=O(\abs{k}^{-m})\)，解析函数则按 \(e^{-a\abs{k}}\) 几何衰减。图上那两条线就是这两句话的直接可视化，蓝线的斜率是 \(-a/\ln 10\)。

这个 \(a\) 是多少，实轴上看不出来。\hyperref[sec:08-Numerical-Analysis/09-Complex-Analysis-and-Residues/01]{``解析延拓与 Fourier 衰减''}会把它算清楚：\(a\) 等于函数解析延拓后最近的那个奇点到实轴的距离，几何收敛的公比就是 \(e^{-a}\)。这里只需记住结论的方向——两个在实轴上一样光滑、一样有界的函数，谱方法上的表现可以差出六七倍的项数，差别写在复平面上。真要提速，除了加 \(N\)，还有一条路是换变量把奇点推远。

\subsection{谱方法是一种赌注}

前面这些拼起来是一句判断：\textbf{谱方法不是更好的差分，它是一种赌注：赌解足够光滑}。赌赢了拿指数收敛，赌输了连二阶差分都比不过。

赌输的样子有两级。解只有有限阶导数时，收敛退化成代数阶，就是图里那条橙线——难看，但还在收敛。解带跳跃时更糟：Fourier 部分和在间断处出现固定比例的过冲，跳跃高度的 \(9\%\) 左右，\(N\) 加多少都不消失，只是把振荡挤得更窄。这就是 Gibbs 现象，它在最大模下根本不收敛。

更麻烦的是误差的位置。全局基意味着一个局部缺陷会污染全局：解在 \(x_0\) 附近有个角点，误差不会乖乖待在 \(x_0\) 附近，而是以振荡的形式铺满整个区间，因为每个基函数都铺满整个区间。有限差分和有限元反过来，局部性差的另一面是局部误差留在局部。所以复杂几何、内部激波、局部尖层、不规则系数这些场合，全局谱方法通常不是合适的工具，得退到有限元、有限体积，或者把区域切开的谱元方法——在每一片上仍然赌光滑，只是把赌注下小了。

这个结构在别的地方也见过。用一组固定的全局基去表示一个对象，收益取决于对象是否真落在这组基擅长的形态里；假设对上了就有几个数量级的便宜可占，对不上时收益直接归零而不是打个折扣。谱卷积一类把算子放到频域里学的网络架构吃的正是这份红利，也承担同样的风险：训练数据里的场若带着尖锐边界，频域截断先把边界磨掉了。

\subsection{非线性会把截断掉的高频送回来}

线性问题里每个模式各走各的，截断很干净。非线性项不然。若 \(u\) 只保留到最高频率 \(K\)，乘积 \(u^2\) 的频率成分最高到 \(2K\)。网格表示不了这些模式，它们不会消失，而是折叠回低频，混进本来算得好好的那些系数里。

这与\hyperref[sec:08-Numerical-Analysis/08-Fourier-and-Spectral-Methods/02]{``采样、混叠与可恢复性''}里的混叠是同一个折叠机制，区别只在于高频的来源：那里是信号本身带的，这里是非线性现场造出来的。加密网格不解决问题，因为网格一密，非线性造出的新频率也跟着往上跑。

处理办法是在做乘积之前把空间加宽：把系数零填充到 \(3N/2\) 长度的网格上，在细网格上做逐点乘法，变回频域后截掉多余部分。等价的说法是 Fourier 方法的 \(2/3\) 规则——只信任低三分之二的模式，剩下的每步清零。谱元和有限元里对应的做法是过积分。代价是每步多两次 FFT，换来的是高频尾部不再无中生有地抬起来。

\subsection{非周期区间换基之后，代价写在条件数和步长上}

解不周期时硬用 Fourier 基，两端拼不上，人为造出一个跳跃，Gibbs 立刻登场。\([-1,1]\) 上的标准替代是 Chebyshev 基，节点向端点聚集，密度按 \(N^{-2}\)。\hyperref[sec:08-Numerical-Analysis/03-Interpolation-and-Approximation/02]{``Runge 现象与节点选择''}给过这么排的理由，那里的结论到这里原样可用：等距节点的高次插值会在端点炸开，Chebyshev 节点不会。

有三处要预先看清。微分矩阵 \(D\) 是稠密的，一次求导 \(O(N^2)\)，不像 Fourier 那样有 FFT 兜底。二阶微分矩阵的条件数按 \(O(N^4)\) 涨，\(N\) 稍大一点，右端的舍入误差就被放大到掩盖谱精度的水平。显式时间推进还会被端点附近那些挤在一起的节点卡住：对流问题的稳定步长是 \(O(N^{-2})\)，扩散问题是 \(O(N^{-4})\)，这与\hyperref[sec:08-Numerical-Analysis/06-ODE/03]{``稳定区域与刚性方程''}里的机制完全相同，最快的那个模态决定所有人的步长。空间上的谱精度不会顺带买到时间上的稳定性，两者要分开算账。

\subsection{系数尾部是最便宜的诊断}

谱方法有一个别的离散化没有的好处：它的中间量本身就是诊断信号。把 \(\abs{\hat u_k}\) 对 \(k\) 画在半对数坐标上，尾部的形状直接告诉你赌注赢没赢。直线下降到某处转平，说明已经吃到舍入极限，加 \(N\) 是浪费；从头就只按幂律缓降，说明解不够光滑；尾部先降后抬，通常是混叠或噪声。

尾部停止衰减的常见原因不止一条，值得逐个排除：

\begin{itemize}
\tightlist
\item
  解本身只有有限阶光滑，或含角点、间断；
\item
  周期边界处函数值或导数拼不上，人为制造了间断；
\item
  非线性项没有去混叠；
\item
  输入数据的噪声达到平台，谱系数只是把它显示了出来；
\item
  舍入误差已经主导，继续加 \(N\) 只会让条件数更差。
\end{itemize}

配套的最低验证也简单：先用单个基函数当输入，检查微分算子是否精确到舍入尺度；再确认边界条件与零频兼容条件；然后拿解析、有限光滑、带跳跃三个测试解各跑一遍，看三条误差曲线是否分别呈现直线、缓降、不降。空间加密和时间步加密必须分开做，否则时间离散的误差会被误记在空间谱方法头上。

至此谱方法的账算完了：它用全局光滑性换来了指数精度，也把全局性写进了自己的每一个系数。而这份全局性还有一层后果没提——系数 \(\hat u_k\) 告诉你频率 \(k\) 上有多少能量，却完全不说这些能量出现在哪一段区间上。对稳定的物理场这不要紧，对一段语音、一条设备振动记录、一列成分随时间变化的信号，这就是致命的。下一节从这个缺口开始。
